\documentclass[10.5pt]{article}
\usepackage{amsmath, amsfonts, amssymb,amsthm}
\usepackage[includeheadfoot]{geometry} % For page dimensions
\usepackage{fancyhdr}

\fancyhf{}
\lhead{Teaching Statement}
\rhead{Tighe McAsey}
\pagestyle{fancy}

% Page dimensions
\geometry{a4paper, margin=1in}
\title{Grappling With Mathematics}
\author{Tighe McAsey}
\date{}

\begin{document}
    \maketitle

    Atleast once a week, I gather around in a circle with ten or twenty other wrestling enthusiasts and watch someone demonstrate a new way to fling or contort my adversaries. This new move is presented in a sequence of roughly four steps and after a couple demos seems crystal clear. Eager to try out a new trick I find my usual partner and jump straight into step one, but usually by the end of what should be step two I'm all tangled up with no clear idea of how I got there. It usually requires a few of the right questions and numerous failed attempts until I can arrive at something vaguely resembling the demo. Later that practice I'll usually try to apply my new knowledge in live wrestling and end up pinned on my back; but thats the only way to learn.

    As a Calculus student I rarely skipped class, and generally I felt I had a crystal clear understanding of what was going on in lecture, because I understood what was going on so clearly I had no issue using Wolfram Alpha on the WebWork (final answer only homework) problems, theres no sense in doing busywork. I failed the midterms.
\end{document}